% Copyright 2026 Sreeram Anil
% SPDX-License-Identifier: CC-BY-4.0
% =============================================================================
%  CHAPTER TEMPLATE — one chapter per estimator. Copy to
%  report/chapters/<family>/<name>.tex and fill every section. Keep the section
%  order; delete only sections that truly do not apply (say so in the text).
%
%  Available notation macros (report/preamble.tex):
%    \vx \vu \vy \vw \vv  bold vectors x u y w v        \mF \mH \mP \mQ \mR \mK \mS  bold matrices
%    \xhat \xminus \xplus = \hat{x}, \hat{x}^-, \hat{x}^+    \Pminus \Pplus
%    \E{.} expectation, \N{\mu}{\Sigma} Gaussian, \T transpose, \inv inverse, \norm{.}
%    \soc = z (state of charge), \ocv = \mathrm{OCV}, \dt = \Delta t
%  Auto-generated result files (do not edit by hand; they exist after the benchmark run):
%    tables/cell_<ESTNAME>.tex        per-scenario metrics table for this estimator
%    figures/cell_<ESTNAME>_nominal.png, figures/cell_<ESTNAME>_init_error.png, ...
%  where <ESTNAME> is the registered estimator name with non-alphanumerics removed.
% =============================================================================
\chapter{<Full algorithm name> (<ACRONYM>)}
\label{ch:<estname>}

\section*{At a glance}
\begin{tabular}{@{}l>{\raggedright\arraybackslash}p{0.76\textwidth}@{}}
Family & <classical Kalman / embedded robust / observer / ...> \\
Header & \texttt{include/estkit/<family>/<file>.hpp} \\
Registered name & \texttt{<ESTNAME>} \\
State dimension & $n_x = 4$ (cell model) — generic in $n_x, n_u, n_y$ \\
Cost per step & $\mathcal{O}(\cdot)$, <number> flops for the cell model, no heap \\
Primary references & \cite{key1,key2} \\
\end{tabular}

\section{Motivation and aerospace/BMS relevance}
Why this method exists, what failure mode of the simpler methods it addresses, and where it is
used (avionics, GNC heritage, automotive BMS ...).  2--4 paragraphs.

\section{Problem statement and assumptions}
State precisely the model class the method assumes (linear/nonlinear, Gaussian/heavy-tailed,
bounded/unbounded uncertainty), and the quantities it estimates.  Define every symbol on first
use.

\section{Derivation}
Full derivation with numbered equations.  Every symbol defined where it first appears
("where $\mK_k\in\mathbb{R}^{n_x\times n_y}$ is the Kalman gain ...").  Include the key lemma /
optimality argument or stability proof sketch from the primary reference, with citation.

\section{Algorithm}
\begin{algorithm}[H]
\caption{<ACRONYM> — one sampling period}
\begin{algorithmic}[1]
\Require $\xplus_{k-1}, \Pplus_{k-1}, \vu_{k-1}, \vy_k, \vu_k$
\State \textbf{Time update:} ...
\State \textbf{Measurement update:} ...
\Ensure $\xplus_k, \Pplus_k$
\end{algorithmic}
\end{algorithm}

\section{Application to the battery model}
How the generic algorithm is instantiated for $\vx=[\soc, v_1, v_2, h]^\T$, $\vu=[i,T]^\T$,
$y = v$ (Chapter~\ref{ch:ecm}); the specific Jacobians/tuning; which \texttt{Options} were used
in the benchmark and why.

\section{Implementation notes}
Numerical safeguards (symmetrisation, Joseph form, Cholesky jitter...), memory footprint,
flop count, float32 behaviour, what happens on an MCU (Cortex-M7 estimate), limitations.

\section{Benchmark results}
\input{tables/cell_<ESTNAME>.tex}
\begin{figure}[H]\centering
\includegraphics[width=\textwidth]{figures/cell_<ESTNAME>_nominal.png}
\caption{<ESTNAME> on the nominal eVTOL mission: true and estimated SOC (top), estimation error with the $\pm 2\sigma$ band when available (middle), voltage residual (bottom).}
\end{figure}
\begin{figure}[H]\centering
\includegraphics[width=\textwidth]{figures/cell_<ESTNAME>_init_error.png}
\caption{<ESTNAME> with a 35\,\% initial SOC error.}
\end{figure}
Discussion of the numbers: which scenarios it handles well, where it fails and why (tie the
explanation back to the assumptions), comparison with the EKF baseline.

\section{Summary}
Three to five sentences: when to use it, when not to.
